\documentclass[11pt]{article}
\usepackage{varnernotes}

\newcommand{\R}{\mathbb{R}}

\begin{document}

\lecturenotes{11a}{Option Sensitivities: Delta, Gamma, Theta, Vega, and Rho}{Fall 2026}{November 3, 2026}{Jeffrey D. Varner}

\learningobjectives
  {Define the Greeks}
  {Interpret delta, gamma, theta, vega, and rho as local derivatives with explicit units and sign conventions.}
  {Approximate option P\&L}
  {Use a second-order Taylor expansion to map small market-factor changes into an estimated premium change.}
  {Compute robust sensitivities}
  {Estimate lattice Greeks with finite differences and diagnose step-size, exercise-boundary, and quote-scaling effects.}

\section{A Premium Is a Surface, Not a Single Number}

An option premium changes when the market state changes. Let
\begin{equation}
  \label{eq:value-function}
  V=V(S,\sigma,\tau,r;K,\delta)
\end{equation}
denote a model value in USD/share, where $S>0$ is spot in USD/share, $\sigma>0$ is annualized volatility, $\tau:=T-t\geq 0$ is time to maturity in years, $r\in\R$ is the continuously compounded risk-free rate, $K>0$ is strike in USD/share, and $\delta\in\R$ is continuous dividend yield. Contract terms $K$ and $T$ are fixed; the displayed market inputs can move.

The Greeks describe the local slope and curvature of the value surface:
\begin{equation}
  \label{eq:greek-definitions}
  \Delta:=\frac{\partial V}{\partial S},
  \qquad
  \Gamma:=\frac{\partial^2V}{\partial S^2},
  \qquad
  \Theta:=\frac{\partial V}{\partial t}=-\frac{\partial V}{\partial\tau},
  \qquad
  \nu:=\frac{\partial V}{\partial\sigma},
  \qquad
  \rho:=\frac{\partial V}{\partial r}.
\end{equation}
Vega is conventionally denoted by the Greek letter $\nu$ even though ``vega'' is not a Greek letter. The definition of theta in \eqnref{greek-definitions} uses \emph{calendar time} $t$: one day passes while expiration $T$ stays fixed, so $\tau$ decreases. Some software instead reports $\partial V/\partial\tau$; its sign is opposite. Always inspect the convention.

\begin{remark}{Raw derivatives versus displayed quotes}{quote-scaling}
If $\sigma$ and $r$ are stored as decimals, raw $\nu$ and $\rho$ measure a full $1.00$ change. A platform that quotes sensitivity per one percentage point reports $\nu/100$ and $\rho/100$. Annual theta is often divided by $365$ or by a trading-day convention. Contract-level dollar exposure also multiplies per-share Greeks by the multiplier, commonly $100$ shares/contract for standard U.S. equity options, and by the signed position count.
\end{remark}

\section{Local P\&L: Slope, Curvature, and Cross-Factor Risk}

For small changes $\Delta S$, $\Delta\sigma$, $\Delta t$, and $\Delta r$, a second-order Taylor approximation that retains spot curvature is
\begin{equation}
  \label{eq:taylor-pnl}
  \Delta V\approx
  \Delta\,\Delta S
  +\frac12\Gamma(\Delta S)^2
  +\Theta\,\Delta t
  +\nu\,\Delta\sigma
  +\rho\,\Delta r.
\end{equation}
This is not an exact identity and is not a first-order total differential because it includes the gamma term. It omits higher-order terms and cross-partials such as vanna, the sensitivity of delta to volatility. The approximation deteriorates for large moves, long holding periods, jumps, or motion across an American exercise boundary.

\begin{figure}[ht]
  \centering
  \begin{tikzpicture}[x=0.95cm,y=0.72cm,>=Latex,font=\sffamily\small]
    \draw[->,vnmuted] (0,0)--(8.1,0) node[right] {$S$};
    \draw[->,vnmuted] (0,0)--(0,5.0) node[above] {$V$};
    \draw[vnlink,very thick]
      plot[smooth] coordinates {(0.4,0.35) (1.5,0.45) (2.6,0.7) (3.7,1.25) (4.8,2.15) (5.9,3.25) (7.2,4.65)};
    \coordinate (p) at (4.8,2.15);
    \fill[vncarnelian] (p) circle (2pt);
    \draw[vncarnelian,thick] (2.7,0.16)--(6.9,4.14);
    \draw[vnmuted,dashed] (4.8,0)--(p);
    \node[below] at (4.8,0) {$S_0$};
    \node[vncarnelian,above left] at (3.5,1.2) {tangent slope $\Delta$};
    \node[vnlink,above left] at (6.8,4.4) {curvature $\Gamma$};
  \end{tikzpicture}
  \caption{Delta is the tangent slope of the premium as a function of spot; gamma measures how that slope changes. A delta-only approximation follows the red tangent, while delta--gamma retains local curvature.}
  \label{fig:delta-gamma}
\end{figure}

For a signed position of $m\in\mathbb Z$ contracts with multiplier $q>0$ shares/contract, the position's Greek vector is $mq$ times the per-share long-option vector. Thus short-position Greeks are the negatives of the corresponding long-position Greeks when both use the same valuation mark and model inputs.

\section{Economic Meaning of the Five Greeks}

\paragraph{Delta $\Delta$.} Delta measures first-order directional exposure. A long vanilla call usually has $0\leq\Delta\leq 1$, while a long vanilla put usually has $-1\leq\Delta\leq 0$ under standard assumptions. A position containing one long option contract and $-q\Delta$ shares is locally delta-neutral: its first-order response to a small spot move is approximately zero. Because delta changes with spot and time, the hedge must be rebalanced.

\paragraph{Gamma $\Gamma$.} Gamma measures the instability of delta. Long vanilla calls and puts generally have nonnegative gamma away from American exercise kinks. A delta-neutral but gamma-positive position benefits from sufficiently large moves before costs and other factor changes; a gamma-negative position is exposed to convex losses and increasingly adverse hedge adjustments.

\paragraph{Theta $\Theta$.} Calendar-time theta measures value change as time passes with the modeled state otherwise fixed. Long optionality commonly has negative theta, reflecting the loss of remaining opportunities, but American exercise, dividends, rates, and deep moneyness can complicate blanket sign claims. Theta is a partial derivative, not a forecast of tomorrow's realized P\&L.

\paragraph{Vega $\nu$.} Vega measures sensitivity to the volatility input used by the pricing model. Long vanilla options generally have positive vega. In market applications, ``implied volatility went up one point'' means $\Delta\sigma=0.01$, not $1.00$. Volatility is itself a surface over strike and maturity, so a single parallel-vega shock is an approximation.

\paragraph{Rho $\rho$.} Rho measures rate sensitivity. European calls typically have positive rho and European puts negative rho under the standard BSM model, but dividends and American exercise can alter magnitude and even local behavior. A single scalar $r$ also compresses an entire discount curve; longer-dated portfolios may require bucketed rate sensitivities.

\begin{center}
\small
\begin{tabular}{llll}
  \hline
  Greek & Perturbation & Raw units & Typical long-vanilla sign \\
  \hline
  $\Delta$ & spot $S$ & share/share & call $+$; put $-$ \\
  $\Gamma$ & spot curvature & share/USD & $+$ \\
  $\Theta$ & calendar time $t$ & USD/share/year & usually $-$ \\
  $\nu$ & volatility $\sigma$ & USD/share per $1.00$ & $+$ \\
  $\rho$ & rate $r$ & USD/share per $1.00$ & call $+$; put $-$ \\
  \hline
\end{tabular}
\end{center}

\begin{remark}{Delta is not generally an exercise probability}{delta-probability}
Delta is a hedge ratio. In the European BSM model with dividend yield $\delta$, call delta is $e^{-\delta\tau}\Phi(d_1)$, whereas the risk-neutral probability of finishing ITM is $\Phi(d_2)$; $d_1\neq d_2$ whenever $\sigma\sqrt\tau>0$. The real-world probability uses another measure, and an American option can exercise before expiration. Delta may serve as a rough market heuristic for ITM probability in limited settings, but it is not the definition of probability of exercise or profit.
\end{remark}

\section{Finite Differences on an American-Option Lattice}

Closed-form Greeks are unavailable for a general American contract, but a consistent pricing function can be perturbed. Let $h_S,h_\sigma,h_t,h_r>0$ be bump sizes in the units of their respective inputs. Central differences give
\begin{align}
  \label{eq:finite-differences}
  \widehat\Delta
  &=\frac{V(S+h_S)-V(S-h_S)}{2h_S},
  &
  \widehat\Gamma
  &=\frac{V(S+h_S)-2V(S)+V(S-h_S)}{h_S^2},\\
  \widehat\nu
  &=\frac{V(\sigma+h_\sigma)-V(\sigma-h_\sigma)}{2h_\sigma},
  &
  \widehat\rho
  &=\frac{V(r+h_r)-V(r-h_r)}{2h_r}.
\end{align}
Arguments not displayed in \eqnref{finite-differences} are held fixed. A one-day calendar theta can be approximated directly by
\begin{equation}
  \label{eq:theta-day}
  \widehat\Theta_{\mathrm{day}}
  :=V(S,\sigma,\tau-1/365,r)-V(S,\sigma,\tau,r),
  \qquad \tau\geq 1/365.
\end{equation}
This quote is a one-day value change, not an annual derivative; multiplying it by $365$ only approximates annual theta when local linearity is adequate.

Each bumped valuation must rebuild the lattice with identical contract conventions and sufficient resolution. A bump that is too large introduces truncation bias; one that is too small amplifies floating-point noise and tree-discretization jumps. Sweep several bump sizes and lattice depths. American values are maxima of exercise and continuation branches, so derivatives can be discontinuous or numerically unstable near the exercise boundary.

\notebooklink
  {L11a: NVDA Call and Put Greeks}
  {https://github.com/varnerlab/CHEME-5660-CourseRepository-Fall-2026/blob/main/lectures/week-11/L11a/CHEME-5660-L11a-Example-Greeks-PutCall-NVDA-Fall-2026.ipynb}
  {Compute delta, gamma, theta, vega, and rho for long and short American call and put positions using CRR repricing.}

\section{Risk Reports Need More Than Five Numbers}

Greeks are model-dependent, point-in-time approximations. A useful report records the valuation timestamp, spot, yield curve, dividend treatment, volatility input, exercise model, contract multiplier, position sign, bump definition, and units. It should show both per-contract and portfolio totals and distinguish model Greeks from broker-supplied quote Greeks.

Stress tests complement local sensitivities. Reprice the portfolio under joint spot--volatility moves, passage of time, rate shifts, and changes in liquidity. Compare the full repricing with \eqnref{taylor-pnl}; the residual reveals higher-order and cross-factor effects. Hedging decisions must also include discrete trading, bid--ask spreads, commissions, jumps, and model error.

\Needspace{0.38\textheight}
\keytakeaways
  {In this lecture, we treated the Greeks as local derivatives of a model-value surface and used them to approximate and diagnose option risk.}
  {Conventions are part of the number}
  {Theta sign, day-count, volatility/rate scaling, multiplier, and position sign must accompany every reported Greek.}
  {Delta--gamma is a local approximation}
  {Delta supplies slope and gamma supplies spot curvature; large moves and exercise-boundary changes require full repricing.}
  {Delta is a hedge ratio, not a probability identity}
  {ITM, exercise, and profit probabilities depend on a distribution, measure, contract policy, and threshold that delta alone does not specify.}
  {Next, combine contract payoffs, profit thresholds, probability models, and delta hedging at the portfolio level.}

\begin{readinglist}
  \item Fischer Black and Myron Scholes, \href{https://doi.org/10.1086/260062}{``The Pricing of Options and Corporate Liabilities,''} \emph{Journal of Political Economy} 81(3), 637--654 (1973).
  \item Robert C. Merton, \href{https://doi.org/10.2307/1831029}{``Theory of Rational Option Pricing,''} \emph{Bell Journal of Economics and Management Science} 4(1), 141--183 (1973).
  \item Paul Wilmott, \emph{Paul Wilmott on Quantitative Finance}, 2nd ed., Wiley (2006), sections on option sensitivities and hedging.
\end{readinglist}

\vfill
{\sffamily\footnotesize\color{vnmuted}\textbf{Educational use and risk notice.} These notes are for instruction only and do not constitute investment advice. Greeks are model-dependent local approximations and options trading involves substantial risk.}

\end{document}
